\chapter{2010年福建卷（文）}

\section{单选题}

\begin{problem}
若集合 \(A = \{x \mid 1 \leqslant x \leqslant 3\}\)，\(B = \{x \mid x > 2\}\)，则 \(A \cap B\) 等于（\quad）

\choices
    {\(\{x \mid 2 < x \leqslant 3\}\)}
    {\(\{x\mid x\geqslant 1\}\)}
    {\(\{x\mid 2\leqslant x < 3\}\)}
    {\(\{x\mid x > 2\}\)}
\end{problem}

\begin{problem}
计算 \(1 - 2\sin^2 22.5^\circ\) 的结果等于（\quad）

\choices
    {\(\frac{1}{2}\)}
    {\(\frac{\sqrt{2}}{2}\)}
    {\(\frac{\sqrt{3}}{3}\)}
    {\(\frac{\sqrt{3}}{2}\)}
\end{problem}

\begin{problem}
若一个底面是正三角形的三棱柱的正视图如图所示，则其侧面积等于（\quad）

\choices
    {\(\sqrt{3}\)}
    {2}
    {\(2\sqrt{3}\)}
    {6}
\end{problem}

\begin{problem}
i 是虚数单位，\( \frac{(1+\i)^{4}}{1-\i} \) 等于（\quad）

\choices
    {i}
    {-i}
    {1}
    {\(-1\)}
\end{problem}

\begin{problem}
若 \(x, y \in \R\)，且 \(\left\{ \begin{array}{l} x \geqslant 1, \\ x - 2y + 3 \geqslant 0, \\ y \geqslant x, \end{array} \right.\) 则 \(z = x + 2y\) 的最小值等于（\quad）

\choices
    {2}
    {3}
    {5}
    {9}
\end{problem}

\begin{problem}
阅读如图所示的程序框图，运行相应的程序，输出的 i 值等于（\quad）

\choices
    {2}
    {3}
    {4}
    {5}
\end{problem}

\begin{problem}
函数 \( f(x) = \begin{cases} x^2 + 2x - 3, & x \leqslant 0, \\ -2 + \ln x, & x > 0, \end{cases} \) 的零点个数为（\quad）

\choices
    {3}
    {2}
    {1}
    {0}
\end{problem}

\begin{problem}
若向量 \(a = (x, 3)\) (\(x \in \R\))，则 \(x^2 = 4\) 是 \(^{17}\) 项 \(a^2 = 5\) 的（\quad）

\choices
    {充分而不必要条件}
    {必要而不充分条件}
    {充要条件}
    {既不充分也不必要条件}
\end{problem}

\begin{problem}
若某校高一年级 8 个班参加合唱比赛的通知分如茎叶图所示，则这组数据的中位数和平均数分别是（\quad）

\choices
    {91.5 和 91.5}
    {91.5 和 92}
    {91 和 91.5}
    {92 和 92}
\end{problem}

\begin{problem}
将函数 \( f(x) = \sin (\omega x + \varphi) \) 的图象向左平移 \( \frac{\pi}{2} \) 个单位，若所得图象与原图象重合，则 \( \omega \) 的值不可能等于（\quad）

\choices
    {4}
    {6}
    {8}
    {12}
\end{problem}

\begin{problem}
若点 O 和点 F 分别为椭圆 \( \frac{x^{2}}{a^{2}} + \frac{y^{2}}{b^{2}} = 1 \) 的中心和左焦点，点 P 为椭圆上的任意一点，则 \( \overrightarrow{OP} \cdot \overrightarrow{FP} \) 的最大值为（\quad）

\choices
    {2}
    {3}
    {6}
    {8}
\end{problem}

\begin{problem}
设非空集合 \( S = \{x \mid m \leqslant x \leqslant l\} \) 满足：当 \( x \in S \) 时，有 \( x^2 \in S \). 给出如下三个命题： \circled{1} 若 \(m = 1\) ，则 \(S = \{1\}\) \circled{2} 若 \(m = -\frac{1}{2}\)，则 \(\frac{1}{4} \leqslant l \leqslant 1\); \circled{3} 若 \(l = \frac{1}{2}\)，则 \(-\frac{\sqrt{2}}{2} \leqslant m \leqslant 0\). 其中正确命题的个数是（\quad）

\choices
    {0}
    {1}
    {2}
    {3}
\end{problem}

\section{填空题}

\begin{problem}
若双曲线 \(\frac{x^2}{4} - \frac{y^2}{b^2} = 1 (b > 0)\) 的渐近线方程为 \(y = \pm \frac{1}{2} x\)，则 \(b\) 等于 \fillinblank{} \fillinblank{}.
\end{problem}

\begin{problem}
将容量为 \(n\) 的样本中的数据分成6组，绘制频率分步直方图. 若第一组至第六组数据的频率之比为 \(2:3:4:6:4:1\)，且前三组数据的频数之和等于 27，则 \(n\) 等于 \fillinblank{} \fillinblank{}.
\end{problem}

\begin{problem}
对于平面上的点集 \(\Omega\)，如果连接 \(D\) 中任意两点的线段必定包含于 \(\Omega\)，则称 \(\Omega\) 为平面上的点集，给出平面上 4 个点集的图形如图所示（阴影区域及其边界），其中为凸集的是 \fillinblank{}. （写出所有凸集相应图形的序号）

\circled{1}

\circled{2}

\circled{3}

\circled{4}
\end{problem}

\begin{problem}
观察下列等式: \circled{1} \( \cos2\alpha=2\cos^{2}\alpha-1; \) \circled{2} \( \cos4\alpha=\cos^{4}\alpha-\cos^{2}\alpha+1; \) \circled{3} \( \cos6\alpha = 32\cos^{6}\alpha - 48\cos^{4}\alpha + 18\cos^{2}\alpha - 1; \) \circled{4} \( \cos8\alpha = 128\cos^{8}\alpha - 256\cos^{6}\alpha + 160\cos^{4}\alpha - 32\cos^{2}\alpha + 1; \) \circled{5} \( \cos10\alpha = m\cos^{10}\alpha - 1280\cos^{8}\alpha + 1120\cos^{6}\alpha + n\cos^{4}\alpha + p\cos^{2}\alpha - 1. \) 可以推测， \( m - n + p = \)  \fillinblank{}.
\end{problem}

\section{解答题}

\begin{problem}
数列 \( \{a_{n}\} \) 中， \( a_{1}=\frac{1}{3} \) ，前 n 项和 \( S_{n} \) 满足 \( S_{n+1}-S_{n}=\left(\frac{1}{3}\right)^{n-1}(n\in\mathbf{N}^{*}) \) . (1) 求数列 \( \{a_{n}\} \) 的通项公式 \( a_{n} \) . 以及前 n 项和 \( S_{6} \) ; (2) 若 \( S_{1}, t(S_{1} + S_{2}), 3(S_{2} + S_{3}) \) 成等差数列，求实数 t 的值.
\end{problem}

\begin{problem}
设平面向量 \(a_{m} = (m,1)\)，\(b_{n} = (2,n)\)，其中 \(m, n \in \{1,2,3,4\}\). (1) 请列出有序数组 \((m, n)\) 的所有可能结果; (2) 记“使得 \( a_{m} \perp (a_{m} - b_{n}) \) 成立的 \( (m, n)^{n} \) 为事件 A，求事件 A 发生的概率.
\end{problem}

\begin{problem}
已知抛物线 \(C: y^{2} = 2px (p > 0)\) 过点 \(A(1, -2)\). (1) 求抛物线 C 的方程，并求其准线方程； (2) 是否存在平行于 OA (O 为坐标原点) 的直线 t，使得直线 l 与抛物线 C 有公共点，且直线 OA 与 l 的距离等于 \( \frac{\sqrt{5}}{5} \) \(\perp\) 若存在，求出直线 l 的方程. 若不存在，说明理由.
\end{problem}

\begin{problem}
某港口 \(O\) 要将一件事要物品用小艇送到一艘正在航行的轮船上. 在小艇出发时，轮船位于港口 \(O\) 北偏西 \(30^{\circ}\) 且与该港口相距 20 海里的 \(A\) 处，并以 30 海里/小时的航行速度沿正东方向匀速行驶. 假设该小艇沿直线方向以 \(v\) 海里/小时的航行速度匀速行驶，经过 \(t\) 小时与轮船相遇. (1) 若希望相遇时小艇的航行距离最小，则小艇航行速度的大小应为多少\(\perp\) (2) 为保证小艇在 30 分钟内 (含 30 分钟) 能与轮船相遇，试确定小艇航行速度的最小值; (3) 是否存在 \(v_{t}\) 使得小艇以 \(v\) 海里/小时 \(v\) 的航行速度行驶. 总能有两种不同的航行方向与轮船相遇\(\perp\) 若存在，试确定 \(v\) 的取值范围; 若不存在. 请说明理由.
\end{problem}

\begin{problem}
已知函数 \( f(x) = \frac{1}{3} x^3 - x^2 + ax + b \) 的图象在点 \( P(0, f(0)) \) 处的切线方程为 \( y = 3x - 2 \). (1) 求实数 \(a, b\) 的值; (2) 设 \( g(x) = f(x) + \frac{m}{x - 1} \) 是 \([2, +\infty)\) 上的增函数. \circled{1} 求实数 m 的最大值; \circled{2} 当 \(m\) 取最大值时，是否存在点 \(Q\)，使得过点 \(Q\) 的直线若能与曲线 \(y = g(x)\) 围成两个封闭图形，则这两个封闭图形的面积总相等\(\perp\) 若存在，求出点 \(Q\) 的坐标; 若不存在，说明理由.
\end{problem}

\begin{problem}
如图，在长方体 \(ABCD - A_{1}B_{1}C_{1}D_{1}\) 中，\(E, H\) 分别是棱 \(A_{1}B_{1}, D_{1}C_{1}\) 上的点 (点 \(E\) 与 \(B_{1}\) 不重合)，且 \(EH // A_{1}D_{1}\). 过 \(EH\) 的平面与棱 \(BB_{1}, CC_{1}\) 相交，交点分别为 \(F, G\). (1) 证明: \(AD //\) 平面 \(EFGH\); (2) 设 \(AB = 2AA_{1} = 2a\). 在长方体 \(ABCD - A_{1}B_{1}C_{1}D_{1}\) 内随机选取一点. 记该点取自几何体 \(A_{1}ABFE - D_{1}DCGH\) 内的概率为 \(p\). 当点 \(E, F\) 分别在棱 \(A_{1}B_{1}, B_{1}B\) 上运动且满足 \(EF = a\) 时，求 \(p\) 的最小值.
\end{problem}

